\providecommand{\rank}{\operatorname{r}}
\providecommand{\RR}{\mathbb R}
\providecommand{\CC}{\mathbb C}
\providecommand{\KK}{\mathbb K}

% ---------- Exercise 71 ----------
\setcounter{enumi}{70}
\item {\bfseries\boldmath
Consider the vector space $\RR^3$. State, with justification, whether each of the following assertions is true or false.
}

\begin{enumerate}[label=\alph*)]

\item $(4,4,4)$ is a linear combination of $\bigl((1,2,1),(1,0,1)\bigr)$.

\begin{answer}
Put $u=(1,2,1)$ and $v=(1,0,1)$. Since
$\alpha u+\beta v=(\alpha+\beta,2\alpha,\alpha+\beta)$,
\[
 P=\langle u,v\rangle=\{(x,y,z)\in\RR^3:x=z\}.
\]
Indeed, for $(x,y,x)$ take $\alpha=y/2$, $\beta=x-y/2$.

True: $(4,4,4)=2u+2v$.
\end{answer}

\item $(1,2,2)$ is a linear combination of $\bigl((1,2,1),(1,0,1)\bigr)$.

\begin{answer}
False: $(1,2,2)$ does not satisfy $x=z$.
\end{answer}

\item $\langle(1,2,1),(1,0,1)\rangle=\{(1,2,1),(1,0,1)\}$.

\begin{answer}
False: $0\in P$ but $0\notin\{u,v\}$, so $P\ne\{u,v\}$.
\end{answer}

\item $\langle(1,2,1),(1,0,1)\rangle\subseteq\{(1,2,1),(1,0,1)\}$.

\begin{answer}
False: the same zero vector disproves $P\subseteq\{u,v\}$.
\end{answer}

\item $\langle(1,2,1),(1,0,1)\rangle\supseteq\{(1,2,1),(1,0,1)\}$.

\begin{answer}
True: $u=1u+0v$ and $v=0u+1v$, so $\{u,v\}\subseteq P$.
\end{answer}

\item $\langle(3,0,-1)\rangle=\{(3,0,-1)\}$.

\begin{answer}
False: $0\in\langle(3,0,-1)\rangle$ but $0\notin\{(3,0,-1)\}$.
\end{answer}

\item $\langle(0,0,0)\rangle=\{(0,0,0)\}$.

\begin{answer}
True: every scalar multiple of $(0,0,0)$ is $(0,0,0)$.
\end{answer}

\item $(4,4,4)\in\langle(1,2,1),(1,0,1)\rangle$.

\begin{answer}
True, by (a).
\end{answer}

\item $(1,2,2)\in\langle(1,2,1),(1,0,1)\rangle$.

\begin{answer}
False, by (b).
\end{answer}

\item $(0,0,0)\in\langle(1,2,1),(1,0,1)\rangle$.

\begin{answer}
True: $0=0u+0v$.
\end{answer}

\item $\langle(4,4,4),(1,2,2)\rangle=\langle(1,2,1),(1,0,1)\rangle$.

\begin{answer}
False. The left span contains $(1,2,2)$, which is not in $P$.
\end{answer}

\item $\langle(1,1,-1),(0,0,0)\rangle=\langle(2,2,-2)\rangle$.

\begin{answer}
True. Adding the zero generator changes no linear combinations,
and multiplication of a generator by $2\ne0$ does not change its span.
\end{answer}

\item $\langle(1,1,-1),(2,2,0)\rangle=\langle(0,0,-2),(3,3,3),(3,3,-1)\rangle$.

\begin{answer}
True. Both sides equal the plane
$Q=\{(x,y,z):x=y\}$. For the left side,
\[
 (x,x,z)=-z(1,1,-1)+\tfrac{x+z}{2}(2,2,0).
\]
For the right side, $(0,0,-2)$ and $(3,3,3)$ already generate $Q$, since
\[
 (x,x,z)=\tfrac{x-z}{2}(0,0,-2)+\tfrac x3(3,3,3).
\]
The additional generator $(3,3,-1)$ also lies in $Q$.
\end{answer}

\item $\langle(1,1,-1),(2,2,0)\rangle=\langle(3,3,3),(3,3,-1),(1,2,-1)\rangle$.

\begin{answer}
False. Every vector in the left span has $x=y$;
the right span contains $(1,2,-1)$, which does not.
\end{answer}

\item $\langle(1,1,1),(1,-1,-1)\rangle=\langle(1,1,-1),(1,-1,1)\rangle$.

\begin{answer}
False. Every vector in the right span satisfies $y=-z$,
but $(1,1,1)$ belongs to the left span and does not satisfy this condition.
\end{answer}
\end{enumerate}
